\documentclass[12pt,a4paper]{article}
\usepackage[T2A]{fontenc}
\usepackage[utf8]{inputenc}
\usepackage[russian]{babel}
\usepackage{geometry}
\usepackage{amsmath,amssymb,mathtools}
\usepackage{array,booktabs,longtable}
\usepackage{xcolor}
\usepackage{tikz}
\usepackage[colorlinks=true,linkcolor=blue!55!black]{hyperref}
\geometry{left=19mm,right=19mm,top=19mm,bottom=21mm}
\setlength{\parindent}{0pt}
\setlength{\parskip}{5pt}
\setlength{\emergencystretch}{2em}
\newcommand{\err}{\textcolor{red}{Ошибка:}}
\newcommand{\idea}{\textcolor{blue!60!black}{Ключевая идея:}}
\newcommand{\outcome}{\textcolor{teal!60!black}{Итог:}}
\newcommand{\subpart}[1]{\par\medskip\textbf{#1}\par\nopagebreak[2]}
\begin{document}
\begin{titlepage}
  \centering
  \vspace*{58mm}
  {\Huge\bfseries Ревью домашней работы\par}
  \vspace{20mm}
  {\Large Автор: Лёвин Артём Александрович\par}
  \vspace{9mm}
  {\large 16 октября 2026 года\par}
  \vfill
  \begin{tikzpicture}
    \draw[blue!45!black,line width=0.7pt] (0,0) .. controls (2,0.8) and (4,-0.8) .. (6,0);
  \end{tikzpicture}
\end{titlepage}
\newpage
\section*{Сводная таблица}
В таблице указаны только задания, в которых найдены ошибки.
\begingroup\small\setlength{\tabcolsep}{4pt}\renewcommand{\arraystretch}{1.35}
\begin{longtable}{@{}p{14mm}p{32mm}p{52mm}p{27mm}p{28mm}@{}}
\toprule
№ & Тема & Суть ошибки или неполноты & Ваш ответ & Правильный ответ\\
\midrule
\endhead
\hyperref[task:7]{7} & Периметр, смешанные числа & Сторона $4\frac16$ заменена на $4\frac12$ & $13\frac18$ см & $12\frac{19}{24}$ см\\
\hyperref[task:9]{9} & Скорость лодки и течение & Для движения против течения скорости сложены вместо вычитания & Для обоих направлений $194\frac54$ м/мин & По течению $195\frac14$, против $179\frac34$ м/мин\\
\bottomrule
\end{longtable}
\endgroup
\clearpage
\thispagestyle{empty}
\null\vfill
\begin{center}{\Huge\bfseries Разбор ошибок}\end{center}
\vfill\null
\clearpage

\phantomsection\label{task:7}
\section*{Задание № 7. Периметр треугольника}
\subpart{Текст задания}
Найдите периметр треугольника со сторонами $4\frac16$ см, $5\frac38$ см и $3\frac14$ см.
\subpart{Ваше решение}
На фотографии записано сложение трёх сторон, но первая длина переписана как $4\frac12$ см:
\[
P=4\frac12+5\frac38+3\frac14
=12+\frac48+\frac38+\frac28
=12+\frac98=13\frac18\text{ см}.
\]
\subpart{Ваш ответ}
$13\frac18$ см.
\subpart{Ошибка}
\err\ При переносе условия в запись периметра сторона $4\frac16$ см заменена на $4\frac12$ см.
\subpart{Где именно возникает ошибка}
Последний правильный шаг --- выбор формулы периметра как суммы трёх длин сторон. Первый неверный переход появляется в самой подстановке значений:
\[
\underbrace{4\frac16+5\frac38+3\frac14}_{\text{длины из условия}}
\quad\longrightarrow\quad
\underbrace{4\frac12+5\frac38+3\frac14}_{\text{Ваша запись}}.
\]
После этой замены Вы последовательно вычисляете сумму уже других чисел.
\subpart{Почему это неверно}
Одна шестая и одна вторая --- разные дроби. Одна вторая равна трём шестым, поэтому подставленная величина больше исходной на две шестых, то есть на одну треть сантиметра:
\[
\frac12-\frac16=\frac36-\frac16=\frac26=\frac13.
\]
Именно на эту величину завышается периметр. Ошибка не в принципе сложения, а в неточном переписывании исходных данных.
\subpart{Ключевая идея}
\idea\ Периметр треугольника равен сумме его трёх сторон. Перед приведением дробей к общему знаменателю нужно проверить, что каждая длина перенесена из условия без изменений.
\subpart{Исправление}
Сохраним выбранный Вами способ и исправим только первый член суммы:
\[
P=4\frac16+5\frac38+3\frac14.
\]
\subpart{Полное правильное решение}
Сложим отдельно целые части:
\[
4+5+3=12.
\]
Для дробных частей знаменатели равны шести, восьми и четырём. Наименьшее общее кратное этих чисел --- двадцать четыре. Поэтому
\[
\frac16=\frac4{24},\qquad
\frac38=\frac9{24},\qquad
\frac14=\frac6{24}.
\]
Сумма дробных частей:
\[
\frac4{24}+\frac9{24}+\frac6{24}
=\frac{4+9+6}{24}=\frac{19}{24}.
\]
Числитель меньше знаменателя, поэтому дополнительная единица не образуется. Получаем
\[
P=12+\frac{19}{24}=12\frac{19}{24}\text{ см}.
\]
Числа девятнадцать и двадцать четыре взаимно просты; дробь сокращать не нужно.
\subpart{Проверка}
В Вашем решении получилось $13\frac18$ см. Проверим разность двух результатов:
\[
13\frac18-12\frac{19}{24}
=13\frac3{24}-12\frac{19}{24}
=12\frac{27}{24}-12\frac{19}{24}
=\frac8{24}=\frac13\text{ см}.
\]
Это в точности разница между ошибочно записанной стороной и стороной из условия, что подтверждает место ошибки.
\subpart{Итог}
\outcome\ Правильный периметр: $12\frac{19}{24}$ см.
\subpart{Задача на закрепление}
Найдите периметр треугольника со сторонами $3\frac15$ см, $4\frac7{12}$ см и $2\frac16$ см.
\clearpage

\phantomsection\label{task:9}
\section*{Задание № 9. Лодка и течение реки}
\subpart{Текст задания}
Моторная лодка в стоячей воде проходит $1500$ м за $8$ мин. Скорость течения реки $7\frac34$ м/мин. Найдите скорости лодки по течению и против течения.
\subpart{Ваше решение}
Вы правильно начали с вычисления скорости лодки в стоячей воде:
\[
\frac{1500}{8}=187\frac12\text{ м/мин}.
\]
Затем для движения по течению сложили скорость лодки и скорость течения, получив запись $194\frac54$ м/мин. Для движения против течения также записано сложение скоростей:
\[
(187+7)+\left(\frac24+\frac34\right)=194\frac54.
\]
Таким образом, два направления в Вашей работе привели к одному числу.
\subpart{Ваш ответ}
По течению --- $194\frac54$ м/мин; против течения --- $194\frac54$ м/мин. Запись смешанного числа с дробной частью больше единицы нестандартна: $194\frac54=195\frac14$.
\subpart{Ошибка}
\err\ Для движения против течения Вы сложили скорости вместо того, чтобы вычесть скорость течения из собственной скорости лодки.
\subpart{Где именно возникает ошибка}
Последний правильный смысловой шаг --- нахождение собственной скорости лодки как расстояния, делённого на время, и сложение скоростей для движения по течению. Первый неверный шаг --- выбор знака плюс при расчёте движения против течения:
\[
\underbrace{187\frac12+7\frac34}_{\text{использовано против течения}}
\quad\text{вместо}\quad
\underbrace{187\frac12-7\frac34}_{\text{нужно против течения}}.
\]
\subpart{Почему это неверно}
В стоячей воде двигатель обеспечивает скорость $187\frac12$ м/мин относительно воды. Если лодка движется в том же направлении, что и вода, течение дополнительно переносит её вперёд: скорости складываются. Если лодка идёт навстречу течению, вода препятствует продвижению относительно берега: скорость течения необходимо вычитать. При положительной скорости течения скорость против течения обязана быть меньше собственной скорости лодки. Ваше значение больше собственной скорости, следовательно, не может быть правильным для движения против течения.
\subpart{Ключевая идея}
\idea\ Собственная скорость лодки равна пройденному в стоячей воде расстоянию, делённому на время. По течению: собственная скорость плюс скорость течения. Против течения: собственная скорость минус скорость течения.
\subpart{Исправление}
Сохраняем верно найденную собственную скорость и меняем знак действия только для движения против течения:
\[
v_{\text{против}}=187\frac12-7\frac34.
\]
\subpart{Полное правильное решение}
Сначала вычисляем скорость лодки в стоячей воде:
\[
v_{\text{собств}}=\frac{1500}{8}=187\frac48=187\frac12\text{ м/мин}.
\]
Для движения по течению прибавляем скорость воды. Приводим одну вторую к четвертям:
\[
\begin{aligned}
v_{\text{по}}&=187\frac12+7\frac34
=187\frac24+7\frac34\\
&=194+\frac54
=195\frac14\text{ м/мин}.
\end{aligned}
\]
Обратите внимание: пять четвертей --- это одна целая и одна четверть, поэтому $194\frac54$ необходимо записать в стандартном виде $195\frac14$.

Для движения против течения вычитаем скорость воды. У дробной части уменьшаемого две четверти, а у вычитаемого --- три четверти, поэтому занимаем одну единицу:
\[
187\frac24=186\frac64.
\]
Теперь можно вычесть по частям:
\[
\begin{aligned}
v_{\text{против}}&=186\frac64-7\frac34\\
&=(186-7)+\frac{6-3}{4}
=179\frac34\text{ м/мин}.
\end{aligned}
\]
Обе скорости положительны; при движении по течению скорость больше собственной, против течения --- меньше.
\subpart{Проверка}
Полусумма скоростей по течению и против течения должна равняться собственной скорости лодки:
\[
\frac{195\frac14+179\frac34}{2}
=\frac{375}{2}=187\frac12\text{ м/мин}.
\]
Половина их разности должна равняться скорости течения:
\[
\frac{195\frac14-179\frac34}{2}
=\frac{15\frac12}{2}=7\frac34\text{ м/мин}.
\]
Оба условия выполнены.
\subpart{Итог}
\outcome\ По течению --- $195\frac14$ м/мин; против течения --- $179\frac34$ м/мин.
\subpart{Задача на закрепление}
Моторная лодка в стоячей воде проходит $900$ м за $6$ мин. Скорость течения реки $5\frac12$ м/мин. Найдите скорости лодки по течению и против течения.

\vfill
\begin{center}\footnotesize Лёвин Артём Александрович эксклюзивно для Ксении Васильченко\end{center}
\end{document}
